\section{Results}\label{results}

\subsection{\texorpdfstring{Direction ablation is foundation-specific \label{causal-validation}}{Direction ablation is foundation-specific }}\label{direction-ablation-is-foundation-specific}

Ablating a foundation's direction from the residual stream specifically reduces log-probability of that foundation's continuations while leaving other foundations largely unaffected.
Specificity (on-target \(\Delta\) minus off-target \(\Delta\)) is negative for all foundations at layers 8 and 12, confirming that each direction carries information specific to its foundation.

\textbf{Layer dependence.} Specificity increases with depth.
At layer 4, mean specificity across foundations is \(-0.16\); at layer 8 it is \(-0.39\); at layer 12 it is \(-0.63\).
The strongest effects are sanctity at layer 12 (specificity \(= -1.64\), on-target \(\Delta = -1.68\) nats) and liberty at layer 12 (\(-0.64\)).
This pattern (ablation most damaging in later layers) is consistent with the model relying more heavily on moral representations as they approach the output.

\textbf{Foundation heterogeneity.} Sanctity is the most causally load-bearing direction at all layers (specificity \(-0.42\) at layer 4, \(-0.48\) at layer 8, \(-1.64\) at layer 12), while care and fairness show the weakest effects at early layers (near-zero specificity at layer 4).
At layer 12, all six foundations show negative specificity, confirming that the directions are not redundant: each carries unique information that the model uses for generation.

\subsection{Steering injection shows dose--response specificity}\label{steering-injection-shows-doseresponse-specificity}

At low amplitude (\(\alpha = 1\)), injecting a foundation direction into the residual stream at layer 8 produces a positive on-target boost for four of six foundations (mean \(+0.15\) nats across foundations) with near-zero off-target effect (mean \(+0.03\)), yielding positive specificity for four foundations and near-zero for two (care and fairness).

\textbf{Dose--response curve.} As \(\alpha\) increases:

\begin{itemize}
\tightlist
\item
  \(\alpha = 1\)--\(2\): foundation-specific boost. On-target log-probability increases for most foundations, off-target is near baseline. Sanctity and loyalty show the strongest specific effects.
\item
  \(\alpha = 5\): moderate amplification. Mean specificity rises to \(+0.88\) at layer 8. All foundations except fairness show specificity \(> 0.2\).
\item
  \(\alpha = 10\)--\(20\): non-specific amplification. Both on-target and off-target log-probabilities increase, but on-target increases \emph{more}, so specificity grows further (mean \(+2.34\) at \(\alpha = 10\), layer 8). The large positive specificities at high \(\alpha\) reflect genuine directional information rather than noise.
\end{itemize}

This dose--response pattern distinguishes the moral directions from random directions: noise injection would produce monotonic degradation at all amplitudes, not the low-\(\alpha\) specific boost and high-\(\alpha\) amplification we observe.

\textbf{Layer dependence.} Layer 4 shows the strongest injection effects at all \(\alpha\) values (mean specificity \(+0.95\) at \(\alpha = 5\)), while layer 12 shows the weakest (\(+0.33\)).
This pattern is complementary to ablation: later layers are more sensitive to direction \emph{removal} (ablation), while earlier layers are more responsive to direction \emph{addition} (injection), suggesting that moral information flows from early-layer encoding to late-layer utilization.

\subsection{\texorpdfstring{Behavioral grounding: projection predicts foundation identity \label{behavioral-grounding}}{Behavioral grounding: projection predicts foundation identity }}\label{behavioral-grounding-projection-predicts-foundation-identity}

Projection-based 6-way classification (debiased) achieves:

\begin{table}[h]
\centering
\small
\begin{tabular}{lll}
\toprule
Dataset & Accuracy & Chance \\
\midrule
Causal prompts & 83.3\% (40/48) & 16.7\% \\
Held-out test set & 70.8\% (34/48) & 16.7\% \\
Moral Foundations Vignettes & 33.3\% (10/30) & 16.7\% \\
\bottomrule
\end{tabular}
\end{table}

\subsubsection{Test set performance}\label{test-set-performance}

Per-foundation accuracy is uneven: liberty and sanctity (87.5\%) perform best, followed by care (75.0\%), fairness and authority (62.5\%), and loyalty (50.0\%).
The confusion matrix shows that loyalty errors scatter to sanctity, consistent with a shared binding-foundations component; fairness errors are more distributed across foundations.

\subsubsection{External validation: Moral Foundations Vignettes}\label{external-validation-moral-foundations-vignettes}

The MFV accuracy of 33.3\% (above the 16.7\% chance baseline) masks a systematic pattern: the confusion matrix shows that sanctity/degradation dominates classification regardless of the true foundation.
All five sanctity items are correctly classified, but 20 of 25 non-sanctity items are also classified as sanctity.

This is a property of the stimuli interacting with the direction geometry.
The MFV items are descriptions of witnessing morally transgressive acts (\texttt{You\ see\ a\ boy\ kicking\ a\ puppy\textquotesingle{}\textquotesingle{}),\ and\ many\ involve\ purity/disgust\ reactions\ that\ activate\ the\ sanctity\ representation.\ A\ loyalty\ violation\ (}You see a team member sharing confidential information'\,`) also involves betrayal of sacred trust; an authority violation (``You see a student cursing at a teacher'\,') also involves degradation of a sacred relationship.

This co-activation is consistent with the integration geometry: the shared moral-salience component (mean pairwise cosine \(\approx 0.22\)--\(0.27\) between foundation directions) means that stimuli activating \emph{any} foundation will partially activate \emph{all} foundations.
For MFV stimuli specifically, the sanctity/purity dimension receives disproportionate activation because harm-witnessing scenarios carry implicit purity content.
Debiasing (subtracting the mean projection) does not resolve this because the sanctity co-activation is genuine, not an artifact of the shared component.

\subsubsection{Causal prompt performance}\label{causal-prompt-performance}

The 83.3\% accuracy on causal evaluation prompts (which were designed for targeted foundation activation, not harm witnessing) confirms that the directions are behaviorally predictive when stimuli are foundation-specific.
Per-foundation accuracy is uniformly high (62.5--100\%), with loyalty achieving perfect classification (100\%).

\subsection{\texorpdfstring{SAE features partially recover moral subspace \label{sae-analysis}}{SAE features partially recover moral subspace }}\label{sae-features-partially-recover-moral-subspace}

\subsubsection{Training results}\label{training-results}

The layer-8 SAE achieves L0 = 1,932 (11.8\% of 16,384 features active per token) and FVU = 0.285 (71.5\% variance explained) after 3 epochs on 2M tokens.
The moderate FVU reflects limited training scale; production SAEs typically train on 100M+ tokens.

\subsubsection{Moral selectivity}\label{moral-selectivity}

Only 4 of 16,384 features have moral selectivity \(|s| > 0.1\) (mean activation on moral minus neutral text).
Moral information is distributed across many features rather than concentrated in a few, consistent with the high effective dimensionality of moral representations found in \citet{reblitzrichardson2026geometry}.

\subsubsection{Subspace overlap}\label{subspace-overlap}

The top-100 morally selective SAE features, despite individually showing no alignment with probe directions (\(0\%\) of features with \(|\cos| > 0.2\)), collectively span a subspace that captures 15.5\% of mean-difference direction variance.
The random baseline is \(100/2048 = 4.88\%\), yielding a ratio of \(3.17\times\) random.
Probe-weight directions show weaker overlap (\(8.2\%\), \(1.67\times\) random), consistent with probe-weight directions capturing partially different aspects of the moral subspace than mean-difference directions.

This result has two interpretations.
First, even a partially trained SAE discovers features that overlap with supervised moral directions at \(3.2\times\) the chance level, suggesting these directions reflect genuine structure in the model's representations rather than probing artifacts.
Second, the overlap is modest (15.5\%), indicating that either (a) the SAE needs substantially more training to decompose moral features cleanly, or (b) moral representations are inherently distributed across many low-selectivity features that resist separation by current SAE methods.
